%& -shell-escape
% !TeX root = thesis-abstract.tex
% !TeX encoding = UTF-8
% !TeX program = XeLaTeX

\documentclass{.local/lib/classes/ndvr-super-article-standalone}

\begin{document}

    \noindent
    В~работе приводятся методы 
    обнаружения аномалий в~видео: 
    \begin{itemize}
        \item {
            простые пороговые методы 
            (рисунок \ref{fig:sad-chart});
        }
        \item {
            методы с адаптивным порогом 
            (рисунок \ref{fig:rescaling-chart});
        }
        \item {
            неконтролируемые статистически методы 
            (рисунки \ref{fig:mean-sgn-diff} 
            и \ref{fig:mean-sgn-diff-fade});
        }
        \item {
            методы на~основе деревьев~регрессии;
        }
        \item {
            методы размножения оценок.
        }
    \end{itemize}

    На~основе проведенных экспериментов~выбран 
    наиболее подходящий по~качеству определения аномалий 
    и~по~вычислительной сложности метод
    поиска на~основе сравнения скользящих 
    средних оценок сигнала.

    Из рисунка~\ref{fig:sad-chart} видно,
    что пороговые методы плохо реагируют на~шум.
    Величину порога требуется вычислить заранее. 
    Величина~порога, которая хорошо определяет аномалии на~видео 
    с~резкой сменой событий, не подходит для~плавного видео. 

    Для методов~c масштабированием по~размаху  
    приходится подбирать размер скользящего окна: 
    при малом размере окна~—~происходит 
    масштабирование паразитных шумов, 
    при большом размере окна~—~пропускаются некоторые аномалии.
    Это видно на~рисунке~\ref{fig:rescaling-chart}
    в~окрестностях пятидесятой секунды.

    % Метод сравнения сигнала со своей моделью показан 
    % на~рисунке~\ref{fig:estimation-check-chart}.
    % Метод неприменим для~разных типов~видео 
    % и чувствителен к случайным всплескам.
    % Это~видно на~рисунке в~промежутке от сороковой 
    % до шестидесятой секунды.

    Метод поиска событий на~основе 
    сравнения скользящих средних оценок сигнала
    лучше обрабатывает случайные всплески
    и при большом размере скользящих окон не чувствителен
    к шумам.
    Но он обладает недостатками, которые присущи всем
    методам со скользящими окнами: при малом размере ловят шум; 
    при большом~—~могут пропустить событие. 
    Например, на~рисунке~\ref{fig:mean-sgn-diff}
    в~окрестностях двадцатой секунды 
    событие не было идентифицировано.
    На рисунке~\ref{fig:mean-sgn-diff-fade}
    показан случай плавного перехода между съёмками.
    Ни~пороговые методы, ни~прочие статистически методы 
    не находят в~этом случае аномалию.

    Методы на~основе деревьев~регрессии 
    и методы размножения оценок обладают очень большим
    набором параметров. 
    Для решения задачи поиска событий в~видео
    они неприменимы. 


    \begin{figure}[htp]
        \includegraphics[width=0.75\linewidth]
        {.local/img/pdf/common-detector/vector-based/luma-extractor/sad-event-chart/Djadja-0-90.pdf}
        \caption{
            Поиск событий наивным пороговым методом
            в~первых 90 секундах видео «Дядя Стёпа~—~милиционер».
            $Y_{t}$~—~средняя яркость кадра в~момент времени $t$.
        }
        \label{fig:sad-chart}
    \vspace{28pt}
        \includegraphics[width=0.75\linewidth]
        {.local/img/pdf/common-detector/vector-based/luma-extractor/rescaling-event-chart/Djadja-0-90.pdf}
        \caption{
            Поиск событий пороговым методом c масштабирование по~размаху
            на~скользящем окне размером 200 секунд
            в~первых 90 секундах видео «Дядя Стёпа~—~милиционер».
            $Y_{t}$~—~средняя яркость кадра в~момент времени $t$.
        }
        \label{fig:rescaling-chart}
    \end{figure}
    \begin{figure}[htp]
        \includegraphics[width=0.75\linewidth]
        {.local/img/pdf/common-detector/vector-based/luma-extractor/mean-sign-diff-event-chart/Djadja-0-90.pdf}
        \caption{
            Поиск событий сравнением скользящих средних оценок сигнала
            в~первых 90 секундах видео «Дядя Стёпа~—~милиционер».
            $Y_{t}$~—~средняя яркость кадра в~момент времени $t$.
        }
        \label{fig:mean-sgn-diff}
    \vspace{28pt}
        \includegraphics[width=0.75\linewidth]
        {.local/img/pdf/common-detector/vector-based/luma-extractor/mean-sign-diff-event-chart/Djadja-90-120.pdf}
        \caption{
            Поиск событий сравнением скользящих средних оценок сигнала
            в~промежутке от 90 до 120 секунд видео «Дядя Стёпа~—~милиционер».
            $Y_{t}$~—~средняя яркость кадра в~момент времени $t$.
        }
        \label{fig:mean-sgn-diff-fade}
    \end{figure}
\end{document}
